% ==========================================================================
% Squelette d'un CV a une colonne, plusieurs variantes depuis une source.
%
% Le candidat ci-dessous est fictif. Il sert a montrer la forme attendue :
% remplacer son contenu, pas sa structure.
%
% Les variantes se posent par des drapeaux. Ici deux exemples :
%   \ifposte   le CV oriente vers un poste en entreprise
%   \ifstage   le CV oriente vers un stage
% Ajouter ou retirer des drapeaux selon les candidatures reellement visees.
% build.sh les bascule et VERIFIE qu'ils ont bascule : un sed qui echoue en
% silence produit sinon un PDF qui compile et qui ment.
% ==========================================================================
\documentclass[10pt,a4paper]{article}

\newif\ifphoto
\phototrue
\newif\ifposte
\postetrue
\newif\ifstage
\stagefalse
\newif\ifen
\enfalse

% ==========================================================================
% Preambule commun a toutes les variantes d'un meme CV.
% Mis en facteur pour qu'une correction de mise en page vaille partout :
% c'est ce qui permet de produire six ou dix PDF sans les diverger.
%
% Drapeaux attendus, poses par le fichier appelant AVANT \input :
%   \ifphoto   photo dans l'en-tete
%   \ifen      corps en anglais (change la langue de babel et la cesure)
%   plus les drapeaux de variante propres au candidat (\ifalt, \ifstage...)
%
% Identite : les quatre commandes ci-dessous sont les SEULES a modifier pour
% changer de candidat. Rien d'autre dans ce fichier ne porte de donnee
% personnelle.
% ==========================================================================

\newcommand{\nomcandidat}{Prenom NOM}
\newcommand{\villecandidat}{Ville, Region}
\newcommand{\telcandidat}{+33 0 00 00 00 00}
\newcommand{\telcandidaturl}{+33000000000}
\newcommand{\mailcandidat}{prenom.nom@exemple.fr}
% Laisser vide une ligne de profil pour la faire disparaitre de l'en-tete.
\newcommand{\githubcandidat}{https://github.com/identifiant}
\newcommand{\linkedincandidat}{https://www.linkedin.com/in/identifiant/}
% Mots-cles du PDF : ils servent la recherche documentaire cote recruteur.
% Ce ne sont pas des mots-cles caches dans la page, ce qui serait une fraude ;
% ce sont les metadonnees normales du fichier.
\newcommand{\motsclescandidat}{}

\ifen
  \usepackage[english]{babel}
\else
  \usepackage[french]{babel}
\fi
\usepackage{fontspec}

% --- polices ---------------------------------------------------------------
% Latin Modern est livre avec toute distribution TeX, donc le CV compile
% partout sans installation. Pour une identite plus contemporaine, remplacer
% par une police installee sur la machine, par exemple :
%   \setmainfont{Source Serif 4} et \newfontfamily\titref{Source Sans 3}
% Garder une serif pour le corps et une sans serif pour les titres : c'est ce
% contraste qui donne au document son air de piece ecrite et non de gabarit.
% Ligatures = NoCommon : le texte reste extractible caractere par caractere,
% ce qui compte pour les logiciels qui liront le PDF.
\setmainfont{lmroman10-regular.otf}[
  BoldFont       = lmroman10-bold.otf,
  ItalicFont     = lmroman10-italic.otf,
  BoldItalicFont = lmroman10-bolditalic.otf,
  Ligatures      = NoCommon
]
\newfontfamily\titref{lmsans10-regular.otf}[BoldFont=lmsans10-bold.otf,
                                            ItalicFont=lmsans10-oblique.otf,
                                            Ligatures=NoCommon]

\usepackage{microtype}
\linespread{1.04}
\usepackage[a4paper,top=1.05cm,bottom=1.0cm,left=1.35cm,right=1.35cm]{geometry}
\usepackage{xcolor}
\usepackage{enumitem}
\usepackage{graphicx}
\usepackage{tikz}
\usepackage{hyperref}
\usepackage{fontawesome5}
\usepackage{needspace}

% --- couleurs --------------------------------------------------------------
% Une seule couleur d'accent, declinee. Deux couleurs vives sur un CV se
% disputent l'attention et aucune ne la gagne.
\definecolor{encre}{HTML}{15202B}   % le texte courant, presque noir
\definecolor{accent}{HTML}{1F4E79}  % titres de section et employeurs
\definecolor{vif}{HTML}{1D68FF}     % liens et puces, uniquement
\definecolor{gris}{HTML}{5B6B84}    % dates, lieux, mentions secondaires
\definecolor{trait}{HTML}{D8DEE8}   % filets

\hypersetup{colorlinks=true, urlcolor=vif, linkcolor=vif,
            pdfauthor={\nomcandidat},
            pdftitle={\nomcandidat},
            pdfkeywords={\motsclescandidat}}

\pagestyle{empty}
\setlength{\parindent}{0pt}
\color{encre}

% --- logos -----------------------------------------------------------------
% \marque{fichier}{hauteur} : logo en ligne, dans un intitule de poste ou de
%                             projet. La hauteur se regle logo par logo, parce
%                             qu'un logo large et un logo carre n'ont pas la
%                             meme presence a hauteur egale.
% \ecusson{fichier}         : logo d'etablissement, dans une boite fixe de
%                             1,6 cm sur 0,72 cm, rapport d'image conserve.
%                             La boite fixe est ce qui aligne la colonne des
%                             formations quelles que soient les proportions.
\newcommand{\marque}[2]{\raisebox{-0.26\height}{\includegraphics[height=#2]{logos/#1}}}
\newcommand{\ecusson}[1]{%
  \makebox[1.65cm][l]{\includegraphics[width=1.6cm,height=0.72cm,keepaspectratio]{logos/#1}}}

% --- structure -------------------------------------------------------------
% \Needspace evite qu'un titre de section ou un intitule de poste reste seul
% en bas de page, separe de ce qu'il annonce.
\renewcommand{\section}[1]{%
  \Needspace{3\baselineskip}
  \vspace{0.42em}
  {\titref\bfseries\footnotesize\color{accent}\MakeUppercase{#1}}\par
  \vspace{0.14em}
  {\color{trait}\hrule height 0.9pt}
  \vspace{0.34em}
}

% \poste{intitule}{employeur}{dates}{lieu}
\newcommand{\poste}[4]{%
  \Needspace{4\baselineskip}
  {\titref\bfseries\small #1}\hfill{\titref\footnotesize\color{gris} #3}\par
  {\titref\small\color{accent} #2}\hfill{\titref\footnotesize\color{gris} #4}\par
  \vspace{0.1em}
}

% Les faits d'une experience. Pas « missions », pas « taches » : des faits.
\newenvironment{faits}{%
  \begin{itemize}[leftmargin=1.0em, itemsep=0.1em, topsep=0.06em, parsep=0pt,
                  label={\color{vif}\scriptsize\raisebox{0.14em}{\textbullet}}]
  \small
}{\end{itemize}\vspace{0.2em}}

% \comp{intitule}{contenu} : une seule colonne, l'intitule ouvre la ligne.
% Un tableau a deux colonnes serait plus joli et moins lisible par les
% logiciels d'analyse : l'ordre des mots y devient imprevisible.
\newcommand{\comp}[2]{%
  {\footnotesize{\titref\bfseries\color{accent}#1 :} #2\par}\vspace{0.2em}
}

% \form{logo}{texte} : une ligne de formation, logo et texte centres.
\newcommand{\form}[2]{%
  \noindent
  \begin{minipage}[c]{0.11\textwidth}\raggedright\ecusson{#1}\end{minipage}%
  \begin{minipage}[c]{0.88\textwidth}\raggedright\footnotesize #2\end{minipage}\par
  \vspace{0.52em}
}

% \entete{titre}{sous-titre}{ligne de disponibilite}
% Le titre est le poste vise, pas le diplome. Le sous-titre porte trois ou
% quatre domaines, separes par des puces : c'est la ligne que l'oeil lit en
% deuxieme et elle doit situer le profil sans phrase.
\newcommand{\entete}[3]{%
\begin{minipage}[t]{\ifphoto 0.82\textwidth \else \textwidth \fi}
\vspace{0pt}
{\titref\fontsize{21}{23}\selectfont\bfseries \nomcandidat}\par
\vspace{0.14em}
{\titref\small\color{accent}\bfseries #1}\par
{\titref\footnotesize\color{gris}#2}\par
\vspace{0.38em}
{\footnotesize
\faMapMarker*\ \villecandidat \quad
\faPhone*\ \href{tel:\telcandidaturl}{\telcandidat} \quad
\faEnvelope\ \href{mailto:\mailcandidat}{\mailcandidat} \quad
\href{\githubcandidat}{\faGithub\ GitHub} \quad
\href{\linkedincandidat}{\faLinkedin\ LinkedIn}\par
\vspace{0.08em}
\faCheckCircle\ {\color{gris}#3}\par}
\end{minipage}%
\ifphoto
\hfill
\begin{minipage}[t]{0.16\textwidth}
\vspace{0pt}\raggedleft
\begin{tikzpicture}
  \clip (0,0) rectangle ++(2.15,2.3);
  \node[inner sep=0] at (1.075,1.15) {\includegraphics[width=2.15cm]{photo.jpg}};
\end{tikzpicture}
\end{minipage}
\fi
\vspace{0.34em}
{\color{trait}\hrule height 1.1pt}
\vspace{0.38em}
}


\begin{document}

% === EN-TETE ==============================================================
% Le titre est le poste vise. Pas « Etudiant en informatique », pas
% « A la recherche d'un emploi » : le recruteur doit lire en premier le nom
% du poste qu'il a ouvert.
\ifstage
\entete{Analyste de données, orientation santé publique}
       {Statistique appliquée \quad\textbullet\quad Python et R \quad\textbullet\quad Données de cohorte \quad\textbullet\quad Restitution aux équipes}
       {\textbf{Recherche un stage de six mois} \quad\textbullet\quad mars – août 2027 \quad\textbullet\quad français, anglais B2}
\else
\entete{Analyste de données, orientation santé publique}
       {Statistique appliquée \quad\textbullet\quad Python et R \quad\textbullet\quad Données de cohorte \quad\textbullet\quad Restitution aux équipes}
       {\textbf{Disponible en CDI} \quad\textbullet\quad préavis d'un mois \quad\textbullet\quad français, anglais B2}
\fi

% === PROFIL ===============================================================
% Cinq a huit lignes, pas davantage. On y repond a trois questions : qui je
% suis professionnellement, ce que je sais faire de verifiable, ce que je
% cherche. Les termes en gras sont ceux qu'un recruteur cherche des yeux.
{\small Diplômée d'un master de santé publique, je travaille les \textbf{données de cohorte} :
nettoyage, analyse et restitution. Au cours des deux dernières années j'ai traité
\textbf{quatre bases de plus de dix mille lignes}, automatisé leur contrôle de qualité en
\textbf{Python} et produit les tableaux de bord utilisés en réunion de service. Je cherche un poste
où l'analyse sert une décision, et non un rapport de plus.\par}

% === EXPERIENCE ===========================================================
\section{Experience professionnelle}

% Modele d'une entree : objectif, moyens, resultat.
% Chaque puce porte au moins un element verifiable : un chiffre, un outil, un
% livrable. Une puce sans rien de verifiable est une puce a reecrire.
\poste{Chargée d'études statistiques \quad\marque{exemple.png}{0.40cm}}
      {Agence régionale de santé \textbullet\ contrat de projet}{janvier 2025 – aujourd'hui}{Lyon}
\begin{faits}
\item Suivi d'une cohorte de \textbf{3\,200 patients} sur dix-huit mois : plan d'analyse
construit avec les médecins, appariement de \textbf{quatre sources} auparavant rapprochées à la
main, \textbf{bilan mensuel produit en deux jours contre onze}.
\item Contrôle de qualité automatisé en \textbf{Python} (pandas, Great Expectations) :
\textbf{trente-deux règles} de cohérence, rapport d'anomalies envoyé chaque lundi aux équipes de
saisie. \textbf{Lignes rejetées ramenées de 7,4 \% à 1,1 \%} en six mois.
\item Tableaux de bord sous \textbf{R Shiny}, utilisés en réunion de service, et formation de
\textbf{six personnes} à leur lecture. Repris par le service sans mon intervention depuis
septembre.
\end{faits}

\poste{Stagiaire, unité d'épidémiologie \quad\marque{exemple.png}{0.40cm}}
      {Centre hospitalier universitaire}{avril – septembre 2024}{Grenoble}
\begin{faits}
\item Analyse rétrospective de \textbf{870 dossiers} anonymisés, sous la responsabilité d'un
praticien hospitalier. Mémoire soutenu, \textbf{note 17/20}, résultats présentés en staff.
\end{faits}

% === PROJETS ==============================================================
% Les projets personnels comptent, a une condition : qu'ils soient decrits
% comme des realisations, avec leur etat reel. Un projet en cours se dit en
% cours. Un depot vide ne se cite pas.
\section{Projets}

\poste{Tableau de bord ouvert des urgences régionales \quad\marque{exemple.png}{0.40cm}}
      {Projet personnel \textbullet\ \textbf{en ligne} \textbullet\ \href{https://exemple.fr}{exemple.fr}}{2025}{}
\begin{faits}
\item Collecte hebdomadaire des données publiques d'affluence, \textbf{Python} et
\textbf{PostgreSQL}, restitution cartographique. \textbf{Quatre-vingts visiteurs mensuels},
code ouvert et documenté.
\end{faits}

% === COMPETENCES ==========================================================
% Une ligne par domaine, l'intitule ouvre la ligne. Pas de barres de niveau :
% personne ne sait ce que valent quatre etoiles sur cinq, et le recruteur le
% sait aussi.
\section{Competences}

\comp{Analyse}{Statistique descriptive et inférentielle, modèles de survie, appariement de sources, plans d'analyse}
\comp{Outils}{Python (pandas, scikit-learn), R (tidyverse, Shiny), SQL, Git}
\comp{Données de santé}{CIM-10, données de cohorte, anonymisation, RGPD appliqué à la recherche}
\comp{Langues}{Français, langue maternelle \quad\textbar\quad Anglais B2}

% === CERTIFICATIONS =======================================================
% Deux groupes separes, et jamais melanges : ce qui est obtenu, ce qui ne
% l'est pas. Une certification en preparation ne porte une date que si
% l'inscription est confirmee.
\section{Certifications}

\comp{Obtenues}{\textbf{Diplôme universitaire de statistique appliquée}, 2024 \quad\textbar\quad \textbf{Bonnes pratiques cliniques}, 2023}
\comp{En préparation}{\textbf{Certificat en épidémiologie de terrain}}

% === FORMATION ============================================================
\section{Formation}

\form{exemple.png}{\textbf{Master de santé publique}, parcours méthodes quantitatives \textbullet\ {\color{accent}Université de Lyon}, 2022 – 2024}
\form{exemple.png}{\textbf{Licence de mathématiques appliquées} \textbullet\ {\color{accent}Université Grenoble Alpes}, 2019 – 2022}

% === PIED =================================================================
% Facultatif. Utile quand il dit quelque chose du candidat ; inutile quand il
% recopie « lecture, voyages, cinema ».
\vspace{0.2em}
{\color{trait}\hrule height 0.9pt}
\vspace{0.3em}
{\footnotesize\color{gris}\textbf{Centres d'intérêt} \quad Course de fond, deux marathons. Bénévole au secours populaire, permanence hebdomadaire.\par}

\end{document}
